\subsection{Basic properties of a function}
\R{injective / surjective / bijective / inverse}{
\begin{rcp}
\item \hd{Injective:} $f(x_1)=f(x_2)\Rightarrow x_1=x_2$. Shortcut: \emph{strictly} monotone \ra\ injective (show $f'>0$ or $f'<0$).
\item \hd{Surjective onto $Y$:} solve $f(x)=y$ for arbitrary $y\in Y$; or show $\operatorname{im}f=Y$ via IVT $+$ the limits at the ends of the domain.
\item \hd{Bijective} = both \ra\ $f^{-1}$ exists. Compute it by solving $y=f(x)$ for $x$ and swapping names.
\item \hd{$f^{-1}$ is continuous} if $f$ is continuous and strictly monotone on an interval; and differentiable with $(f^{-1})'(y_0)=\frac1{f'(x_0)}$ provided $f'(x_0)\ne0$.
\end{rcp}
Also: \hd{even} $f(-x)=f(x)$, \hd{odd} $f(-x)=-f(x)$; \hd{periodic} $f(x+T)=f(x)$; \hd{bounded} $|f|\le C$.
}
\R{"determine $a,b$ so that $f$ is continuous / differentiable"}{
\begin{rcp}
\item Continuity at the seam $x_0$: set $\lim_{x\to x_0^-}f=\lim_{x\to x_0^+}f=f(x_0)$ \ra\ one equation.
\item Differentiability: additionally $\lim_{x\to x_0^-}f'=\lim_{x\to x_0^+}f'$ \ra\ second equation.
\item Solve the (linear) system for $a,b$. Two unknowns need exactly these two conditions.
\end{rcp}
}

\subsection{Limits of functions}
\D{$\eps$--$\delta$ and one-sided}{
$\lim_{x\to x_0}f(x)=L$ :$\iff$ $\forall\eps>0\ \exists\delta>0:\ 0<|x-x_0|<\delta\Rightarrow|f(x)-L|<\eps$.\\
Limit exists $\iff$ $\lim_{x\to x_0^-}f=\lim_{x\to x_0^+}f$ (and both finite).
$\lim_{x\to\infty}f=L$: $\forall\eps\ \exists K:\ x>K\Rightarrow|f(x)-L|<\eps$.
}
\R{compute $\lim_{x\to x_0}f(x)$}{
\begin{rcp}
\item Plug in $x_0$. No problem \ra\ done (continuity).
\item $\frac00$: factor \& cancel, expand roots by the conjugate, or \ra\ Taylor (\S5.6) / l'Hospital (\S5.5).
\item $\frac{c\ne0}{0}$: $\to\pm\infty$; determine the sign from the left/right \ra\ answer only exists one-sided.
\item $x\to\pm\infty$ rational: divide by highest power in the denominator. $\sqrt{\cdot}$ for $x\to-\infty$: $\sqrt{x^2}=|x|=-x$. \warn
\item $|x|$, $\sgn$, piecewise, floor: \textbf{always} check both one-sided limits.
\item Bounded $\times$ null: squeeze, e.g. $x\sin\frac1x\to0$ but $\sin\frac1x$ has no limit.
\end{rcp}
}

\subsection{Continuity}
\D{continuity + sequential criterion}{
$f$ continuous at $x_0$ :$\iff$ $\forall\eps>0\ \exists\delta>0:\ |x-x_0|<\delta\Rightarrow|f(x)-f(x_0)|<\eps$
$\iff$ $\lim_{x\to x_0}f(x)=f(x_0)$
$\iff$ \hd{(Heine)} for every sequence $x_n\to x_0$ one has $f(x_n)\to f(x_0)$.
}
\R{prove / disprove continuity at $x_0$}{
\hd{Prove:} built from $+,-,\times,\div$ (den.\,$\ne0$), $\circ$ of continuous standard functions \ra\ continuous, one line. Otherwise $\eps$--$\delta$: bound $|f(x)-f(x_0)|\le C|x-x_0|$ and take $\delta=\eps/C$.\\
\hd{Disprove:} find \emph{two} sequences $x_n\to x_0$, $y_n\to x_0$ with $\lim f(x_n)\ne\lim f(y_n)$ (Heine). Classic: $x_n=\frac1{2\pi n}$ vs.\ $y_n=\frac1{\pi/2+2\pi n}$ for $\sin\frac1x$; rationals vs.\ irrationals for Dirichlet.
}
\F{types of discontinuity (Unstetigkeitsstellen)}{
\hd{removable} \de{hebbar}: both one-sided limits exist and are equal $\ne f(x_0)$.
\hd{jump} \de{Sprung}: both exist, differ (this is what monotone functions have; at most countably many).
\hd{essential/oscillating}: at least one one-sided limit fails to exist ($\sin\frac1x$, $\frac1x$).
}

\subsection{The three big theorems on $[a,b]$ --- and when to use which}
\F{cite these verbatim}{
$f:[a,b]\to\Rl$ \textbf{continuous}, $[a,b]$ compact:
\begin{bul}
\item \hd{IVT \de{Zwischenwertsatz}:} $c$ between $f(a)$ and $f(b)$ \ra\ $\exists x\in[a,b]:f(x)=c$. \emph{Use for: existence of a zero / solution / fixed point.}
\item \hd{Extreme value thm \de{Satz vom Max/Min, Weierstrass}:} $f$ is bounded and \emph{attains} $\max$ and $\min$. \emph{Use for: "the maximum is attained", bounding.}
\item \hd{Heine:} $f$ continuous on a \textbf{compact} interval \ra\ $f$ is \emph{uniformly} continuous. \emph{Use to get uniform continuity for free.}
\end{bul}
\warn All three fail on open/unbounded intervals ($f(x)=1/x$ on $(0,1]$).
}
\R{show "$f(x)=g(x)$ has a solution" / "$f$ has a zero"}{
\begin{rcp}
\item Set $h:=f-g$ (fixed point $f(x)=x$: $h(x):=f(x)-x$). State: $h$ continuous as difference of continuous functions.
\item Find two points with opposite signs: evaluate at convenient points, or take limits $x\to\pm\infty$ and argue with the sign for large $|x|$.
\item IVT \ra\ $\exists x^*$ with $h(x^*)=0$. \hd{Existence done.}
\item \hd{Uniqueness:} show $h$ is strictly monotone ($h'>0$ or $h'<0$ everywhere) \ra\ at most one zero. Alternative: Rolle --- two zeros would force $h'(\xi)=0$, contradiction.
\end{rcp}
}
\E{$e^x=p(x)$ --- typical multiple choice}{
\begin{bul}
\item IVT gives a solution \textbf{not always}: $p(x)=-1$ has none ($e^x>0$).
\item $p$ of \emph{odd} degree \ra\ always at least one solution ($e^x-p(x)$ changes sign).
\item Number of solutions is always \emph{finite}? \textbf{Yes} --- $e^x-p(x)$ is analytic, $\ne0$, so finitely many zeros; and for any $N$ there \emph{is} a polynomial with $\ge N$ solutions (take the $N$-th Taylor polynomial of $e^x$-like constructions / oscillating differences).
\end{bul}
}

\subsection{Uniform continuity \de{gleichm\"assige Stetigkeit}}
\D{difference to ordinary continuity}{
\hd{continuous:} $\forall x_0\forall\eps\exists\delta(\eps,x_0)$. \quad
\hd{uniformly continuous:} $\forall\eps\ \exists\delta(\eps)\ \forall x,y:\ |x-y|<\delta\Rightarrow|f(x)-f(y)|<\eps$.\\
\textbf{The $\delta$ must not depend on the point.} unif.\ cont.\ \ra\ cont., never the other way.
}
\R{"is $f$ uniformly continuous?"}{
\begin{rcp}
\item \hd{Domain compact $[a,b]$ and $f$ continuous?} \ra\ YES by Heine. Done, one line.
\item \hd{Is $f'$ bounded on the domain, $|f'|\le L$?} \ra\ MVT gives $|f(x)-f(y)|\le L|x-y|$ (Lipschitz) \ra\ YES, take $\delta=\eps/L$. \emph{This is the workhorse on $\Rl$.}
\item \hd{$f$ continuous on $\Rl$ with finite limits at $\pm\infty$?} \ra\ YES (compactness argument on a large interval + almost constant outside).
\item \hd{Suspect NO?} Find $x_n,y_n$ with $|x_n-y_n|\to0$ but $|f(x_n)-f(y_n)|\ge\eps_0>0$. Classics: $x^2$ on $\Rl$ ($x_n=n$, $y_n=n+\frac1n$); $\frac1x$ on $(0,1)$; $\sin\frac1x$ on $(0,1)$.
\item \hd{Case distinction} if a parameter can degenerate (e.g.\ $n=0$ vs.\ $n\ge1$) --- treat the degenerate case separately.
\end{rcp}
\warn Unbounded $f'$ does NOT prove "not uniformly continuous": $\sqrt x$ on $[0,1]$ has $f'\to\infty$ but is unif.\ cont.\ (compact).
}
\E{$f_n(x)=\dfrac{x}{1+nx^2}$ on $\Rl$: uniformly continuous?}{
$n=0$: $f_0(x)=x$, Lipschitz with $L=1$ \ra\ unif.\ cont.\\
$n\ge1$: $f_n'(x)=\dfrac{1-nx^2}{(1+nx^2)^2}$, so $|f_n'(x)|\le\dfrac{1+nx^2}{(1+nx^2)^2}=\dfrac1{1+nx^2}\le1$.
Bounded derivative \ra\ Lipschitz with $L=1$ \ra\ \textbf{uniformly continuous for every $n$.}
}

\E{disproving uniform continuity: $f(x)=x^2$ on $\Rl$}{
Take $x_n=n$, $y_n=n+\frac1n$. Then $|x_n-y_n|=\frac1n\to0$, but
$|f(x_n)-f(y_n)|=\left|n^2-\left(n+\tfrac1n\right)^2\right|=2+\tfrac1{n^2}\ge2=:\eps_0$.
So no $\delta$ can work for $\eps=2$ \ra\ \textbf{not uniformly continuous}.
(Same pattern for $\frac1x$ on $(0,1)$: $x_n=\frac1n$, $y_n=\frac1{2n}$ \ra\ difference $n$.)
}

\subsection{Sequences of functions $(f_n)$ --- pointwise limit}
\R{find the pointwise limit $f(x)=\lim_n f_n(x)$}{
\begin{rcp}
\item \textbf{Freeze $x$} (treat it as a constant), let $n\to\infty$ using \S2.2.
\item Do a \hd{case distinction in $x$} whenever the answer changes: typically $x=0$ vs.\ $x\ne0$, $|x|<1$ vs.\ $=1$ vs.\ $>1$.
\item Write the limit function piecewise. \hd{If $f$ is discontinuous, convergence cannot be uniform} (see \S4.7).
\end{rcp}
}
\E{pointwise limits you should recognise}{
$\dfrac{x}{1+nx^2}$: $x=0$ \ra\ $0$; $x\ne0$ \ra\ $\frac{x}{nx^2}\to0$. So $f\equiv0$ on $\Rl$.\\
$x^n$ on $[0,1]$ \ra\ $f=0$ on $[0,1)$, $f(1)=1$ (discontinuous \ra\ not uniform).\\
$\frac{\sin(nx)}{n}\to0$; $\ n x e^{-nx^2}\to0$; $\ \arctan(nx)\to\frac\pi2\sgn(x)$.
}

\subsection{Uniform convergence $f_n\rightrightarrows f$}
\D{the only formula you need}{
$f_n\to f$ \hd{uniformly on $D$} :$\iff$
$\displaystyle \|f_n-f\|_\infty:=\sup_{x\in D}|f_n(x)-f(x)|\xrightarrow[n\to\infty]{}0$
$\iff$ $\forall\eps\exists N\forall n\ge N\ \forall x\in D:|f_n(x)-f(x)|<\eps$ (one $N$ for \emph{all} $x$).
}
\R{decide uniform convergence (the standard 3-point task)}{
\begin{rcp}
\item Get the pointwise limit $f$ first (\S4.6). Uniform convergence can only be towards it.
\item Form $g_n(x):=|f_n(x)-f(x)|$ and compute $M_n:=\sup_x g_n(x)$: differentiate $g_n$ w.r.t.\ \textbf{$x$} (n fixed), set $g_n'(x)=0$, plug the critical point back in. Also check the boundary/$\pm\infty$.
\item $M_n\to0$ \ra\ \textbf{uniform}. $M_n\not\to0$ \ra\ \textbf{only pointwise}.
\item \hd{Shortcut for NO:} exhibit one sequence $x_n$ with $|f_n(x_n)-f(x_n)|\ge\eps_0>0$ (usually $x_n$ = the maximiser). One such sequence is a complete disproof.
\item \hd{Shortcut for NO:} $f_n$ all continuous but $f$ is not \ra\ not uniform (contrapositive of the Uniform Limit Theorem).
\end{rcp}
}
\E{$f_n(x)=\frac{x}{1+nx^2}$ on $\Rl$: uniform?}{
$f\equiv0$, so $M_n=\sup_x\left|\frac{x}{1+nx^2}\right|$.
$f_n'(x)=\frac{1-nx^2}{(1+nx^2)^2}=0\Rightarrow x=\pm\frac1{\sqrt n}$.
$f_n\!\left(\frac1{\sqrt n}\right)=\frac{1/\sqrt n}{2}=\frac1{2\sqrt n}$.
So $M_n=\frac1{2\sqrt n}\to0$ \ra\ \textbf{converges uniformly to $0$ on all of $\Rl$.}
}
\E{$f_n(x)=nxe^{-nx^2}$ on $[0,1]$: not uniform}{
Pointwise $f\equiv0$. Maximiser: $x_n=\frac1{\sqrt{2n}}$, $f_n(x_n)=\sqrt{\frac n2}e^{-1/2}\to\infty\ne0$ \ra\ \textbf{only pointwise}.
}
\F{what uniform convergence buys you (and what it does not)}{
$f_n$ continuous on $D$, $f_n\rightrightarrows f$ \ra
\begin{bul}
\item \hd{$f$ is continuous} (Uniform Limit Theorem). \emph{Pointwise alone gives nothing --- $x^n$ on $[0,1]$.}
\item On $[a,b]$: $f$ is Riemann-integrable and $\int_a^b f_n\to\int_a^b f$ (swap $\lim$ and $\int$). \emph{Pointwise alone: false.}
\item \warn \textbf{Derivatives do NOT pass to the limit:} $\frac{\sin(nx)}{n}\rightrightarrows0$ but $f_n'=\cos(nx)$ diverges. You need $f_n'\rightrightarrows g$ \emph{uniformly} plus convergence at one point; then $f'=g$.
\item $\sup|f_n|\ge\frac1{2n}$ does \emph{not} block uniform convergence to $0$ --- $\frac1{2n}\to0$. Only a lower bound that does \emph{not} vanish blocks it.
\item $\max$ attained at a different point $x_n$ each time is irrelevant; only the \emph{value} $M_n$ counts.
\end{bul}
}
\F{Weierstrass $M$-test (for $\sum f_n$)}{
$|f_n(x)|\le M_n\ \forall x\in D$ and $\sum M_n<\infty$ \ra\ $\sum f_n$ converges \textbf{absolutely and uniformly} on $D$. Use for function series / power series on $[-r,r]$.
}
